\section[Notebook 1989-2: the perturbation lemma]{Notebook 1989-2: the perturbation lemma, algebraic form\protect\verifie{Notebook1989_2}}\label{sec:quillen-1989-2}

\emph{Context.} On pages~76--80 of \carnet{1989}{2} Quillen transfers a perturbation $\theta$ of a differential along a homotopy equivalence $(u, v, h)$, with $vu - 1 = [d, h]$, to $\theta' = u(1 - \theta h)^{-1}\theta v$: the basic perturbation lemma.

\begin{encadre}
\begin{theoreme}[after Quillen, notebook 1989-2, pp.~76--80]\label{thm:quillen-1989-2}
In a ring, let $d$, $d_M$, $\delta$, $h$, $i$, $p$ satisfy $d^2 = d_M^2 = 0$, $(d + \delta)^2 = 0$, $di = id_M$, $pd = d_Mp$, $ip = 1 + dh + hd$, and let $1 - \delta h$ be invertible; put $X = (1 - \delta h)^{-1}\delta$. Then $dX + Xd + X^2 + X(dh + hd)X = 0$, and with $d_M' = d_M + pXi$, $i' = i + hXi$, $p' = p + pXh$, $h' = h + hXh$:
\[ d_M'^2 = 0,\quad (d + \delta)i' = i'd_M',\quad p'(d + \delta) = d_M'p',\quad i'p' = 1 + (d + \delta)h' + h'(d + \delta). \]
If moreover $pi = 1$, $hi = 0$, $ph = 0$ and $h^2 = 0$, then $p'i' = 1$.
\end{theoreme}
\end{encadre}

\begin{proof}
$(1 - \delta h)X = \delta$, and $1 - h\delta$ is invertible with inverse $1 + h(1 - \delta h)^{-1}\delta$, with $X(1 - h\delta) = \delta$. Multiplying the first identity by $1 - \delta h$ on the left and $1 - h\delta$ on the right gives $d\delta + \delta d + \delta^2 = (d + \delta)^2 - d^2 = 0$. The four identities reduce to it: each difference is $p(\cdot)i$, $h(\cdot)i$, $p(\cdot)h$ or $h(\cdot)h$ of the first, plus multiples of $(1 - \delta h)X - \delta$ and $X(1 - h\delta) - \delta$.
\end{proof}

\emph{Verification in Lean.} The file \href{\leanblobmain/Notebook1989_2.lean}{\texttt{Notebook1989\_2.\allowbreak{}lean}} proves \texttt{key}, \texttt{sq\_eq\_zero}, \texttt{chain\_left}, \texttt{chain\_right}, \texttt{homotopy} and \texttt{retract}. Everything is stated in one ring (the endomorphisms of the sum of the two complexes), with operators composed without signs; the side conditions are used only for $p'i' = 1$. The page's $\theta'$ is $pXi$. The axioms used are the three usual axioms of mathlib.

\litterature{Classical (Brown, 1965; Gugenheim, 1972; Huebschmann--Kadeishvili, 1991); the edition does not present it as new.}
